\chapter{Results}
This chapter presents the empirical findings and performance metrics gathered during the experimental evaluation of the integrated robotic vision system. To verify the efficacy of the proposed distributed ROS 2 framework, benchmarks were collected across three primary operational axes: the computational efficiency of the perception pipeline, the metric precision of the spatial localization node, and the trajectory reliability of the person-following control loop. The subsequent sections provide a detailed analysis of these quantitative and qualitative results based on the testing procedures established in the previous chapter.

\section{Object Detection Performance}
The performance of the primary perception pipeline closely aligned with the algorithmic expectations of the well-established YOLO framework. However, the practical evaluation highlighted significant hardware-dependent performance variations. Initial implementation attempts focused on executing the object detection models locally on the TurtleBot 4's embedded Raspberry Pi 4 Single-Board Computer. Experimental benchmarks revealed that the onboard ARM-based processor lacked the required computational capacity; even when deploying the highly optimized "nano" ($n$) scale model, the system suffered from critical frame-rate drops and severe latency, rendering closed-loop robotic tracking impossible.

Following the paradigm shift toward a distributed network offloading architecture, the perception workload was transferred to the external workstation. Upon migrating the inference tasks to the workstation's high-performance graphics hardware, the processing capability experienced a substantial upgrade. The system not only achieved fluid real-time frame rates but successfully scaled up to utilize the "extra-large" ($x$) YOLO model configurations. This transition significantly enhanced human feature extraction and classification precision without bottlenecking the robotic middleware.

During continuous tracking sequences under standard ambient lighting, the inference loop operated with high speed and reliability. A minor qualitative edge-case was identified during testing: when the human target was briefly obstructed by tight geometric architecture or household obstacles, the system occasionally registered transient false positives or split detections for a fraction of a second. Aside from these brief latency fluctuations during partial occlusions, the workstation-hosted YOLO pipeline provided a highly stable and dependable semantic stream, establishing a robust foundation for continuous spatial localization.

\section{Localization Accuracy}
The performance of the \texttt{object\_locator} node demonstrated high efficacy in transforming two-dimensional camera bounding boxes and synchronized stereo depth matrices into three-dimensional spatial coordinates. Empirical testing confirmed that the Euclidean distance calculations along the camera's optical axis were highly accurate and remained well within an acceptable operational margin of error for indoor target tracking. 

However, the experimental evaluation revealed a critical architectural challenge regarding real-time data throughput and network latency. Because the localization node simultaneously subscribes to multiple high-bandwidth data streams—including two separate, uncompressed image topics for RGB frames and stereo depth maps—the system is forced to ingest and process a massive volume of data per second. Over continuous operational periods, a noticeable processing lag accumulates within the pipeline, causing the localization coordinates to drift behind the real-time position of the target.

To mitigate this latency, several software optimization strategies were implemented and evaluated:
\begin{itemize}
	\item \textbf{Subscription Management:} Unnecessary or redundant topic subscriptions were programmatically terminated to streamline the node's execution loop.
	\item \textbf{Quality of Service (QoS) Tuning:} The ROS~2 QoS profiles were adjusted to prioritize volatile, best-effort data delivery over historical reliability, allowing the system to drop old frames in favor of the most recent sensor updates.
	\item \textbf{Visualization Deactivation:} Computational overhead was further reduced by disabling all optional image annotation and debug visualization publishers (\texttt{visual = False}).
\end{itemize}

Despite these optimizations, a persistent communication bottleneck remained. Quantitative analysis indicates that the primary constraint is not the CPU processing limits of the workstation, but rather the physical bandwidth limitations of the DDS middleware when transmitting heavy image messages over a wireless network link. The high frequency and sheer size of the arriving image packets saturate the wireless network, introducing transport delays. This transmission lag directly impacts the performance of the downstream \texttt{person\_follower} state machine, as it receives delayed positional updates, thereby challenging the fluid responsiveness of the robot's real-time trajectory control.

\subsection{Person Following Performance}
The empirical evaluation of the centralized Finite State Machine (FSM) demonstrated that after resolving initial configuration challenges, the high-level control logic transitions smoothly and predictably between its five operational states. The logic successfully managed the execution of individual state routines under nominal conditions. However, the practical deployment of the autonomous person-following behavior was heavily impacted by a prominent system bottleneck.

The primary operational constraint stems directly from the previously identified localization latency. Because the 3D coordinate updates arrive at the FSM with a persistent processing lag, the robot's real-time actuation suffers from an algorithmic delay. This latency introduces a severe control loop error during the \texttt{State.SEARCHING} phase. As the TurtleBot~4 executes a continuous rotation in place to scan for a human target, the delayed perception data causes the platform to overshoot the target's actual physical position. By the time the workstation processes the message confirming a human detection, the physical orientation of the robot has already rotated past the camera's field of view. Consequently, this systemic overshoot results in an immediate loss of the target, forcing the FSM back into the searching state and creating an unstable, oscillating behavior.



